\appendix
\section*{Appendix G: Operator-Algebraic Formalization of the Altitude System}

\subsection*{G.0 Overview}

This appendix presents the operator-algebraic formalization of the
altitude architecture used in the Constellation–Manifold Toolkit.
The purpose is to express:
\begin{itemize}
    \item altitudes as operator classes,
    \item legal transitions as operator homomorphisms,
    \item purity as algebraic closure,
    \item non-generativity as identity-only extension,
    \item boundedness as operator-norm constraints,
    \item diagnostics as bounded perturbations,
    \item drift semigroup as a strongly continuous operator semigroup,
    \item invariant extraction as a projection operator,
    \item steering as an automorphism,
    \item collision invariant as a fixed-point operator,
    \item the Series V firewall as an admissible operator subalgebra.
\end{itemize}

\subsection*{G.1 Operator Algebra}

Let $X$ be a Banach space representing the system state space.

Define the operator algebra:
\[
\mathcal{B}(X) = \{ T : X \to X \mid T \text{ bounded linear operator} \}.
\]

Each altitude corresponds to a distinguished operator class:
\[
\mathcal{O}(A_0), \mathcal{O}(A_1), \dots, \mathcal{O}(I_1)
\subseteq \mathcal{B}(X).
\]

\subsection*{G.2 Altitude Operators}

Altitude $A_0$ (Reversible Geometry):
\[
G \in \mathcal{O}(A_0),\quad G^{-1} \in \mathcal{O}(A_0).
\]

Altitude $A_1$ (Conceptual Metric):
\[
M \in \mathcal{O}(A_1),\quad M(x,y) \in \mathbb{R}_{\ge 0}.
\]

Altitude $\Phi_7$ (Curvature Envelope):
\[
E \in \mathcal{O}(\Phi_7),\quad E(x) \in \mathbb{R}.
\]

Diagnostics $D_1, D_2, D_3$:
\[
C,S,\Sigma \in \mathcal{O}(D_i).
\]

Altitude $\Phi_8$ (Drift Semigroup):
\[
D_t \in \mathcal{O}(\Phi_8),\quad D_{t+s}=D_t D_s.
\]

Altitude $\Phi_9$ (Invariant Extraction):
\[
I \in \mathcal{O}(\Phi_9),\quad I(x) \in \mathbb{R}^k.
\]

Altitude $S_1$ (Steering):
\[
\Theta \in \mathcal{O}(S_1),\quad \Theta^{-1} \in \mathcal{O}(S_1).
\]

Altitude $I_1$ (Collision Invariant):
\[
\Omega \in \mathcal{O}(I_1),\quad \Omega(x) \in \{0,1\}.
\]

\subsection*{G.3 Legal Transition Homomorphisms}

Define the altitude transition homomorphism:
\[
\phi_{a,b} : \mathcal{O}(a) \to \mathcal{O}(b)
\]
such that:
\[
\phi_{a,b}(T) = T' \quad \text{if } T \text{ is altitude-safe}.
\]

Allowed transitions:
\[
\phi_{A_0,A_1},\;
\phi_{A_1,\Phi_7},\;
\phi_{\Phi_7,D_1},\;
\phi_{D_1,D_2},\;
\phi_{D_2,D_3},
\]
\[
\phi_{D_3,\Phi_8},\;
\phi_{\Phi_8,\Phi_9},\;
\phi_{\Phi_9,S_1},\;
\phi_{S_1,I_1}.
\]

Forbidden transitions:
\[
\phi_{a,b} \text{ undefined}.
\]

\subsection*{G.4 Purity as Algebraic Closure}

Purity constraint:
\[
P(a) \iff \mathcal{O}(a) \text{ is closed under composition}.
\]

Thus:
\[
T_1,T_2 \in \mathcal{O}(a) \implies T_1 T_2 \in \mathcal{O}(a).
\]

\subsection*{G.5 Non-Generativity as Identity-Only Extension}

Non-generativity requires:
\[
\mathcal{O}(a) \subseteq \mathcal{B}(X)
\quad \text{and} \quad
\mathcal{O}(a) \text{ contains no generative operators}.
\]

Formally:
\[
\mathrm{span}(\mathcal{O}(a)) = \mathrm{span}(\{\mathrm{id}_X\}).
\]

\subsection*{G.6 Boundedness as Operator-Norm Constraint}

Boundedness requires:
\[
\|T\| < C_a \quad \forall T \in \mathcal{O}(a).
\]

Thus altitude operators are uniformly bounded.

\subsection*{G.7 Diagnostics as Bounded Perturbations}

Diagnostics are bounded perturbations:
\[
D_i = \mathrm{id}_X + \Delta_i,
\]
where:
\[
\|\Delta_i\| < \epsilon.
\]

\subsection*{G.8 Drift Semigroup as Strongly Continuous Semigroup}

Define:
\[
D_t \in \mathcal{O}(\Phi_8)
\]
such that:
\[
D_{t+s}=D_t D_s,
\]
and:
\[
\lim_{t\to 0} \|D_t x - x\| = 0.
\]

Thus $\{D_t\}$ is a strongly continuous semigroup.

\subsection*{G.9 Invariant Extraction as Projection Operator}

Invariant extraction is a projection:
\[
I : X \to \mathbb{R}^k,
\]
with:
\[
I^2 = I.
\]

\subsection*{G.10 Steering as Automorphism}

Steering is an automorphism of the operator algebra:
\[
\Theta : \mathcal{B}(X) \to \mathcal{B}(X),
\]
with:
\[
\Theta^{-1} \circ \Theta = \mathrm{id}.
\]

\subsection*{G.11 Collision Invariant as Fixed-Point Operator}

Collision invariant:
\[
\Omega : X \to \{0,1\}
\]
is a fixed-point operator:
\[
\Omega(\Omega(x)) = \Omega(x).
\]

\subsection*{G.12 Firewall as Admissible Operator Subalgebra}

Define the firewall subalgebra:
\[
\mathcal{O}_{\mathrm{safe}} =
\{ T \in \mathcal{B}(X) \mid T \text{ satisfies altitude constraints} \}.
\]

Thus:
\[
\mathcal{O}_{\mathrm{safe}} \subseteq \mathcal{B}(X)
\]
is an admissible operator subalgebra.

\subsection*{G.13 Final Safety Theorem}

\textbf{Theorem.}  
A state $x$ is altitude-safe if and only if:
\[
\Omega(x)=1.
\]

Equivalently:
\[
x \in \mathcal{O}_{\mathrm{safe}}.
\]

This completes the operator-algebraic formalization.
